\section{Preference Data：人类偏好本身就是一个噪声测量系统}

SFT 只能模仿 demonstrations，但人们未必能写出自己最偏好的回答。这就是 generation-value gap：容易评价一个候选是否更好，却不一定容易从零生成最佳答案。Preference learning 把监督改成 $(x,y^+,y^-)$，其中 $y^+$ 优于 $y^-$。

\singleslide{slide-31.jpg}{InstructGPT 流水线进入第二阶段：从 SFT imitation 转向基于反馈的 policy optimization。}{31}

Slide 31 在视觉上重复课程主轴，但教学作用不同：前半讲已经完成 demonstrations，现在开始追问“如果人类写不出最佳答案，能否只比较候选”。这一步把模型从条件语言模型重新解释为 policy；训练目标不再只是复现一份参考文本，而是提高被评价系统判为更好的回答概率。

\slidepair{slide-32.jpg}{slide-33.jpg}{Imitation 与 reward optimization 的目标差异，以及人们会评价却未必会生成的 generation-value gap。}{32--33}

Slide 32 用两个概率目标完成概念切换：SFT 逼近一个可采样的 reference distribution，RLHF 则在 policy 自己的输出分布下最大化 reward。Slide 33 的新闻摘要例子说明为什么需要这种切换：标注者可能很难从空白页写出优秀摘要，却能稳定指出两个候选中哪一个更准确、简洁。比较监督降低了生产标签的认知负担，但也把质量定义交给了 pair 与 rubric。

\begin{importantbox}{Generation-value gap 的含义}
“能评价”不等于“评价无噪声”，只表示比较两个候选通常比生成理想答案容易。Preference pipeline 的价值来自把开放式生成任务改写成局部排序问题；它的风险则是候选集之外的优秀行为不会被看见，而标注者可能用长度、语气等捷径完成比较。
\end{importantbox}

\slidepair{slide-34.jpg}{slide-35.jpg}{RLHF 的课程路线图，以及 preference data 章节要解决的标签类型与质量问题。}{34--35}

Slides 34--35 把后半讲分成三层：先研究 feedback 如何收集，再比较 PPO 与 DPO 如何消费反馈，最后检查过度优化的副作用。这里的顺序不能反过来，因为算法只会放大已有测量；如果 pairwise labels 主要反映 annotator style 或 compensation pressure，换一个更强 optimizer 只会更快地学到这些偏差。

\subsection{Pairwise feedback 的标准形式}

本节先定义 measurement process。给定 prompt $x$，从策略采样多个 responses，annotator 选择更优者。可用 Bradley--Terry model 表示偏好概率：

$$
p(y^+\succ y^-\mid x)
=\sigma\!\left(r_\phi(x,y^+)-r_\phi(x,y^-)\right).
$$

$r_\phi$ 是 reward model，$\sigma$ 是 sigmoid。训练 reward model 时最大化观察到的 pairwise preferences。

\begin{figure}[H]
\centering
\includegraphics[width=0.82\textwidth]{slide-36.jpg}
\caption{标准 pairwise feedback setup：对同一 prompt 比较候选回答。}
\figsource{CS336 Spring 2026 Lecture 15，Slide 36。}
\end{figure}

\begin{knowledgebox}{读图：偏好标签不是客观真值}
标签由 guideline、annotator expertise、time budget、compensation 和 interface 共同生成。不同 annotator 可能优化 factuality、politeness、helpfulness 或安全性中的不同组合，因此 preference data 应被视为带噪测量，而不是绝对 reward。
\end{knowledgebox}

\slidepair{slide-37.jpg}{slide-38.jpg}{InstructGPT 与 Bard 的历史标注指南：pairwise label 背后其实是一套复杂 rubric。}{37--38}

这两页提醒我们，$y^+\succ y^-$ 不是点击按钮时自然产生的真值。InstructGPT guideline 要求同时权衡 helpfulness、truthfulness 与 harmlessness，Bard 的旧标注文档也把问题拆成事实、相关性、风格和安全等维度。不同维度可能冲突：更完整的回答可能更长，更谨慎的回答可能显得不够直接，安全拒绝也可能损害表面 helpfulness。

\begin{warningbox}{Rubric 冲突必须显式处理}
若界面只允许选一个赢家，标注者会在多个维度间做隐式加权，reward model 随后把这种个人权重当作统一目标。更稳健的流程会保存维度级标签、无法判断选项与 adjudication 记录，并对事实正确性使用工具或专家复核，而不是强迫 generalist 猜测。
\end{warningbox}

\slidepair{slide-39.jpg}{slide-40.jpg}{现代数据工人的平台分布，以及 generalist 与 expert annotation 之间巨大的报酬差异。}{39--40}

Slides 39--40 把“human feedback”还原成劳动系统。平台工、承包商和领域专家面对不同任务、培训、计价与保密限制；同一条医学、法律或代码回答，对 generalist 来说可能只能按流畅度判断，对专家来说才可能核验实质正确性。报酬差异并非旁枝，它会决定谁愿意花时间查证、标注者流失率以及数据能否覆盖高难度领域。

\slidepair{slide-41.jpg}{slide-42.jpg}{Crowdsourcing 的质量控制难题，以及大规模数据收集不可忽略的伦理成本。}{41--42}

Slide 41 列出三个直接质量风险：难以验证 annotator 能力、难以确保其真正检查 correctness，以及无法排除其使用 AI 代做。Slide 42 进一步要求把伦理纳入系统边界，包括低报酬、心理伤害、隐私暴露和全球劳动关系。高质量 preference data 不能只靠事后过滤；任务设计、支持机制、申诉与公平报酬本身会影响标签可靠性。

\begin{knowledgebox}{课堂提示：annotation quality 是系统变量}
把 crowd worker 当作可互换的 API 会丢失最重要的 provenance。实际工程中应记录招募渠道、资格测试、领域、地区、培训版本、单题耗时、报酬、agreement 与复核结果；这些 metadata 既用于伦理审计，也用于解释模型为什么在某些用户群体或专业任务上系统偏移。
\end{knowledgebox}

\subsection{Demographics 与 annotator style 会进入模型}

为了进一步说明 measurement bias，课程展示了 annotator population 对模型行为的影响。即使收集很多标签，如果人群结构单一，模型也会把这一群体的价值判断平均化为默认行为。

\begin{figure}[H]
\centering
\includegraphics[width=0.84\textwidth]{slide-43.jpg}
\caption{Annotator demographics 会系统性改变 RLHF 后的行为分布。}
\figsource{CS336 Spring 2026 Lecture 15，Slide 43。}
\end{figure}

Slide 43 的关键不是证明某个人群“更正确”，而是展示聚合标签会把 annotator distribution 写入模型。若训练人群在政治态度、文化背景、专业能力或风险容忍上不代表部署用户，RLHF 后的默认行为就会系统偏向该群体；增加标签数量只能降低采样方差，不能消除 population bias。

\singleslide{slide-44.jpg}{即使标注者数量增加，不同 annotator 的 style preference 仍会显著改变模型排序。}{44}

Slide 44 把 demographics 推进到更细的 style 层。并非所有标注者都偏好相同的长度、礼貌程度、引用方式或解释结构；某些标注者的稳定偏好会在聚合后占主导。因而“many annotators”不自动等于“neutral reward”，应报告分组 agreement、少数偏好和无法达成一致的任务，而不是只保留多数票。

\begin{warningbox}{多数偏好不等于普适偏好}
聚合 preference 时要保存 annotator metadata、agreement、uncertainty 和 slice metrics。单一 reward model 往往压平合法差异；对于价值冲突问题，personalization 或 explicit policy 可能比“平均人类偏好”更诚实。
\end{warningbox}

\subsection{AI Feedback：便宜、稳定，但可能同源偏置}

本节接着看 LM-generated preferences。强模型在 system-level ranking 上可以接近 human agreement，因此被广泛用于 UltraFeedback、Tulu 等 pipeline。它显著降低成本，也能统一 rubric；但 judge 与 policy 可能共享训练偏差，产生 verbosity bias、self-preference 和 style matching。

\begin{figure}[H]
\centering
\includegraphics[width=0.82\textwidth]{slide-45.jpg}
\caption{GPT-4 等模型作为 pairwise judge 时与人类排名的相关性。}
\figsource{CS336 Spring 2026 Lecture 15，Slide 45。}
\end{figure}

Slide 45 给出了使用 AI judge 的现实依据：在 system-level ranking 上，强模型与人类排序可有很高相关性，单题 agreement 也接近 human inter-annotator level。但 system-level correlation 会掩盖特定 prompt 的共同错误；judge 与人类都可能偏好流畅、长且格式漂亮的答案，因此必须保留可验证任务和人工审计集作为独立锚点。

\slidepair{slide-46.jpg}{slide-47.jpg}{UltraFeedback、Tulu 等 AI-feedback pipeline，以及 Constitutional AI 的 self-training 循环。}{46--47}

Slide 46 表明 AI feedback 已不是边缘实验，而是许多开源模型获取大规模 pairwise labels 的基础设施。Slide 47 则展示 self-training 的另一种形式：先用 constitution 或原则生成 critique 与 revision，再把改进后的回答用于训练。两者共同把昂贵的人类判断转成可扩展的模型调用，但也可能在 generator、judge 与 policy 之间形成同源偏置闭环。

\begin{knowledgebox}{老师强调：即使资源充足，也常会使用 AI feedback}
课堂在约 1:02:01 明确指出，现实中的 preference-data 项目即使拥有较强资源，也常因成本、速度与专业覆盖而依赖 AI feedback。正确问题不是“human 还是 AI 二选一”，而是哪些维度由 verifier 检查、哪些由专家审计、哪些可由 judge 扩展，以及三者发生冲突时谁拥有最终裁决权。
\end{knowledgebox}

\singleslide{slide-48.jpg}{RLHF 中显著的长度效应：优化偏好会把已有的 verbosity bias 进一步放大。}{48}

这页把 SFT 章节的 length confounder 带回 RLHF。若 pairwise data 中较长回答更常胜出，reward model 会把长度当作便宜特征，policy optimization 随后继续增加 token 数。验收时应画 win rate 对长度差的条件曲线，并检查“删去冗余后是否仍胜出”；否则 reward 上升可能只是模型更熟练地迎合 judge。

\begin{importantbox}{Preference data 的审计字段}
至少记录 prompt source、candidate policy/version、sampling parameters、annotator/judge、rubric、pair order、response length、agreement 和 adjudication。缺少这些字段时，很难解释 reward model 学到了内容还是格式。
\end{importantbox}

\begin{knowledgebox}{Nota de clase}
课程并不把 human feedback 浪漫化：crowdsourcing 很难验证 factual correctness，expert annotation 成本高，AI use 还会污染“人类标签”。RLHF 的数据问题与 pretraining data 一样，是系统工程和治理问题。
\end{knowledgebox}

\singleslide{slide-49.jpg}{从高质量 pairwise pipeline 进入算法问题：PPO 与 DPO 如何把比较信号写回 policy。}{49}

Slide 49 是数据与优化之间的交接点。到这里我们拥有的仍不是“真实效用函数”，而是一批受候选分布、rubric、人口结构和 judge bias 影响的 pairs。PPO 会把它们拟合成 reward 并在线采样，DPO 会直接在静态 pairs 上更新；两条路线的共同上限仍由数据测量过程决定。

\subsection{本章小结}

Preference data 把“生成答案”转成“比较答案”，但比较仍受 rubric、人群、成本和 judge bias 影响。算法只能优化收到的信号，无法自动修正错误的 measurement process。

